% Copyright (C) 2014-2024 by Thomas Auzinger <thomas@auzinger.name>

\documentclass[draft,final]{vutinfth} % Remove option 'final' to obtain debug information.

% Define convenience functions to use the author name and the thesis title in the PDF document properties.
\newcommand{\authorname}{Alexander Rinsche} % The author name without titles.
\newcommand{\thesistitle}{Context-sensitive Semantic Search in Graphs} % The title of the thesis. The English version should be used, if it exists.

% Create the XMP metadata file for the creation of PDF/A compatible documents.
\begin{filecontents*}[overwrite]{\jobname.xmpdata}
\Author{\authorname}                                    % The author's name in the document properties.
\Title{\thesistitle}                                    % The document's title in the document properties.
\Language{de-AT}                                        % The document's language in the document properties. Select 'en-US', 'en-GB', or 'de-AT'.
\Keywords{a\sep list\sep of\sep keywords}               % The document's keywords in the document properties (separated by '\sep ').
\Publisher{TU Wien}                                     % The document's publisher in the document properties.
\Subject{Thesis}                                        % The document's subject in the document properties.
\end{filecontents*}

% Load packages to allow in- and output of non-ASCII characters.
\usepackage{lmodern}        % Use an extension of the original Computer Modern font to minimize the use of bitmapped letters.
\usepackage[T1]{fontenc}    % Determines font encoding of the output. Font packages have to be included before this line.
\usepackage[utf8]{inputenc} % Determines encoding of the input. All input files have to use UTF8 encoding.

% Extended LaTeX functionality is enables by including packages with \usepackage{...}.
\usepackage{amsmath}    % Extended typesetting of mathematical expression.
\usepackage{amssymb}    % Provides a multitude of mathematical symbols.
\usepackage{mathtools}  % Further extensions of mathematical typesetting.
\usepackage{microtype}  % Small-scale typographic enhancements.
\usepackage[inline]{enumitem} % User control over the layout of lists (itemize, enumerate, description).
\usepackage{multirow}   % Allows table elements to span several rows.
\usepackage{booktabs}   % Improves the typesetting of tables.
\usepackage{subcaption} % Allows the use of subfigures and enables their referencing.
\usepackage[ruled,linesnumbered,algochapter]{algorithm2e} % Enables the writing of pseudo code.
\usepackage[dvipsnames,table]{xcolor} % Allows the definition and use of colors. This package has to be included before tikz.
\usepackage{nag}        % Issues warnings when best practices in writing LaTeX documents are violated.
\usepackage{todonotes}  % Provides tooltip-like todo notes.
\usepackage{morewrites} % Increases the number of external files that can be used.
\usepackage[a-2b,mathxmp]{pdfx}      % Enables PDF/A compliance. Loads the package hyperref and has to be included second to last.
\usepackage[acronym,toc]{glossaries} % Enables the generation of glossaries and lists of acronyms. This package has to be included last.

% Set PDF document properties
\hypersetup{
    pdfpagelayout   = TwoPageRight,           % How the document is shown in PDF viewers (optional).
    linkbordercolor = {Melon},                % The color of the borders of boxes around hyperlinks (optional).
}

\setpnumwidth{2.5em}        % Avoid overfull hboxes in the table of contents (see memoir manual).
\setsecnumdepth{subsection} % Enumerate subsections.

\nonzeroparskip             % Create space between paragraphs (optional).
\setlength{\parindent}{0pt} % Remove paragraph indentation (optional).

\makeindex      % Use an optional index.
\makeglossaries % Use an optional glossary.
%\glstocfalse   % Remove the glossaries from the table of contents.

% Set persons with 4 arguments:
%  {title before name}{name}{title after name}{gender}
%  where both titles are optional (i.e. can be given as empty brackets {}).
\setauthor{}{\authorname}{}{male}
\setadvisor{Prof.}{Emanuel Sallinger}{}{male}

% For bachelor and master theses:
\setfirstassistant{Dr.}{Eleonora Laurenza}{}{female}
%\setsecondassistant{Pretitle}{Forename Surname}{Posttitle}{male}
%\setthirdassistant{Pretitle}{Forename Surname}{Posttitle}{male}

% For dissertations:
%\setfirstreviewer{Pretitle}{Forename Surname}{Posttitle}{male}
%\setsecondreviewer{Pretitle}{Forename Surname}{Posttitle}{male}

% For dissertations at the PhD School and optionally for dissertations:
\setsecondadvisor{Pretitle}{Forename Surname}{Posttitle}{male} % Comment to remove.

% Required data.
\setregnumber{12120519}
\setdate{01}{01}{2025} % Set date with 3 arguments: {day}{month}{year}.
\settitle{\thesistitle}{Kontext-sensitive Semantische Suche in Graphen} % Sets English and German version of the title (both can be English or German). If your title contains commas, enclose it with additional curvy brackets (i.e., {{your title}}) or define it as a macro as done with \thesistitle.
\setsubtitle{Optional Subtitle of the Thesis}{Optionaler Untertitel der Arbeit} % Sets English and German version of the subtitle (both can be English or German).

% Select the thesis type: bachelor / master / doctor.
% Bachelor:
\setthesis{bachelor}
%
% Master:
%\setthesis{master}
%\setmasterdegree{dipl.} % dipl. / rer.nat. / rer.soc.oec. / master
%
% Doctor:
%\setthesis{doctor}
%\setdoctordegree{rer.soc.oec.}% rer.nat. / techn. / rer.soc.oec.

% For bachelor and master:
\setcurriculum{Software \& Information Engineering}{Software \& Information Engineering} % Sets the English and German name of the curriculum.

% Optional reviewer data:
\setfirstreviewerdata{Affiliation, Country}
\setsecondreviewerdata{Affiliation, Country}


\begin{document}

\frontmatter % Switches to roman numbering.
% The structure of the thesis has to conform to the guidelines at
%  https://informatics.tuwien.ac.at/study-services

\addtitlepage{naustrian} % German title page.
\addtitlepage{english} % English title page.
\addstatementpage

\begin{danksagung*}
\todo{Ihr Text hier.}
\end{danksagung*}

\begin{acknowledgements*}
\todo{Enter your text here.}
\end{acknowledgements*}

\begin{kurzfassung}
\todo{Ihr Text hier.}
\end{kurzfassung}

\begin{abstract}
\todo{Enter your text here.}
\end{abstract}

% Select the language of the thesis, e.g., english or naustrian.
\selectlanguage{english}

% Add a table of contents (toc).
\tableofcontents % Starred version, i.e., \tableofcontents*, removes the self-entry.

% Switch to arabic numbering and start the enumeration of chapters in the table of content.
\mainmatter

\chapter{Introduction}
What is the definition of Information? To further specify the question, such a definition should cover the following commonly observed properties of information:
\begin{itemize}
    \item The same sentence contains a different amount of information depending on the receiver and their context
    \item It's verb ''to inform'' means to convey a fact that should influence someone's behavior.
    \item An ''informed decision'' means to be aware of as many relevant facts as possible concerning a certain decision.
\end{itemize} 
By those properties, the definition of information for this thesis will be ''facts or knowledge relevant to a receiver that influence their behavior or decisions''. The amount of information can be determined by the severity of change in their actions depending on getting or missing a said piece of information. For a machine learning model, this can essentially be measured by loss. A data sample with a high loss obviously contains more information than the last ten data samples that the model predicted correctly and therefore had little to no loss. 

From this there is a clear distinction between a Knowledge Base and a Information System: While a Knowledge Base is a collection of facts or, more generally speaking, knowledge, an Information System needs to actively filter the facts relevant to a user or a query. 
%So although a database can be considered a perfectly suitable Information System to experienced developers, for a common user there is too much (to them) uninteresting information and overhead attached to properly use it. 
For a Knowledge Base to become an Information System, a mechanism that produces the relevant information to their users at the moment where it becomes relevant, without additional overhead of filtering out unnecessary data, is paramount. In this thesis we will explore and evaluate one possible mechanism leveraging Text Embeddings and Knowledge Graph Technology.

\section{Motivation}
The recent success of transformer architectures has affected most areas of Computer Science to some degree. In the field of Information Retrieval and Search through document corpi, the technique of Semantic Search has gained new traction. While previous approaches like one hot encodings or continuous bag of words all had their shortcomings, new transformer architectures are able to capture complex relationships and meanings of words in dense vectors and through the attention mechanism that the transformer architectures all share even ambiguities of words are representable in these models.
The price is the high initial cost to train these models, the billions of parameters that need to be loaded in active memory and additional finetuning that is often required when it comes to more specialized domains.
The computational resources required for these tasks are still infeasible to most individuals and even many companies, which is why big tech companies like Google and OpenAI are currently so successful. The aim of this thesis is to propose an alternative architecture where semantic embeddings don't have to be perfect but can be improved later on without fine-tuning of the embedding model itself and rather through exploiting additional information about the domain. The proposed approach should also produce a decent recommendation system and maybe even a semantic search that can take user-based relevance into account. To achieve this, the problem won't be a loose collection of documents, but a set of documents and users organized and linked in a rich taxonomy. By leveraging the information encoded in a neighborhood in such a graph, additional information about the context should be includable through basic Graph Neural Networks (GNNs).


\section{Data Sources}
The necessary data for training and evaluating the model was obtained from qibri.
qibri is a Knowhow Management System targeted at medium to large companies with highly specialized domain-specific knowledge. Its purpose is to manage and unify all documentation of Knowhow in a central place, to create a better overview of how well the core processes and products are documented and provide the necessary documentation to an employee when they needed. It has several mechanisms to encourage users to regularly check and update the Knowhow that they are responsible for, a versioning and approval process system to make sure changes are transparent, can be backtracked and are formally accepted by decision makers. Knowledge retrieval is particularly difficult in this context, however - there's a lot of information that might seem relevant when searching for a specific subject, but usually a lot of it just isn't relevant to the user who started the search. So, although a document might seem to contain a lot of keywords relevant to a given search, it might still be the case that it just isn't relevant due to it's particular scope and target audience within a company. To mitigate this, all documents and users are organized in a taxonomy, containing locations, departments, relevant topics, etc. Using this taxonomy, one can quickly get an idea which documents are more closely related to a user's daily work and which might be less relevant, just by using the distances between them. Unifying a semantic search with the relevance scoring particular to the user triggering the search will be another area of exploration of this thesis, although the necessary data to evaluate such a joint search is currently not available and will only be obtainable through a feedback loop.

\section{Problem Statement}
Finding, optimizing and verifying such a system unfortunately exceeded the scope of a bachelor thesis. Therefore, the focus of this work will be on creating a sensible framework for quickly iterating over different versions. The main challenges this thesis posed were
\begin{description}
    \item[history preservation] i.e. incorporating some kind of versioning so that different extraction runs can easily be compared to each other.
    \item[extendability] building the pipeline in a way such that changes to the data extraction and the embedding process are ''easy'' to implement.
    \item[ease of deployment] - due to the classic difficulties of evaluating a recommender system or a search engine (difficulties defining ''correct'' results, relevances for different people, etc.), one of the main requirements imposed by qibri was the option to quickly deploy to production and to test it there.
\end{description}
In addition to that, there are three iterations of the system that are going to be explicitly compared: A naive Semantic Search, A static Message Passing Contextualization, where no parameters are actually learned and context is statically aggregated and a proper GNN-based approach, where actual learnable parameters are fit. These steps should give a decent first idea if the proposed methodology seems promising or not.


\chapter{Theoretical Background}
This section will focus on establishing the necessary terminology and background information necessary to understand the Machine Learning and AI aspects of this thesis. Firstly, a brief overview over methods to represent text as a low-dimensional vector will be provided. Next, the notion of a Knowledge Graph will briefly be introduced, followed by a detailed description over Graph Neural Networks and how they synergize with text embeddings and a Knowledge Graph.

\section{Text Embeddings}
% Word2Vec
Representing words or, more generally speaking, text, as (relatively) low-dimensional embedding vectors (instead of the trivial method of one-hot-encoding a dictionary) such that relative distances and orientation capture some notion of meaning is one of if not the main topic of Natural Language Processing (NLP). One of the first big successes on this field was Word2Vec, which was published in a paper by Mikolov et al., at Google, in 2013. This technique uses a sliding window approach where every word in the vocabulary of the model gets projected to a dense vector and these vectors are then fit using one of two methods: \textit{Continuous bag of words} or \textit{continuous skip-gram}. Both of these methods use simple linear transformations where the one-hot-encoded word vector gets projected to a dense vector space with dimensionality $\mathcal{D}$, while the bag-of-words approach uses the surrounding to predict a missing word, the skip-gram version uses a specific input word to predict the surrounding. The novel part of these approaches is that the predicitons are computed by another simple linear layer followed by a softmax, without any non-linearities or any larger depth. \cite{word2vec_preprint} This however suffices to learn relationships like $v('France') - v('Paris') + v('Madrid') \approx v('Spain')$, while still being quite efficient to compute. \cite{word2vec_official} The limitations arise from more complex relationships and a deeper understanding of text. By construction, Word2Vec has no mechanism to alter meaning by context, it doesn't even capture the sequence of words. As the term "Bag of Words" implies, the model considers no notion of ordering, words that appear before and after are summed up, averaged, and then the most similar word is predicted. The skip-gram approach also only predicts words that are also likely to appear given a specific input word, not the \textit{next} or \textit{previous} words. In addition to that, words that have several different meanings like ''bear'' or ''bow'' also aren't distinguished. The semantic nuances are accumulated in a single embedding vector and any predictions are ambiguous if they were a result of the verb or the noun. A final drawback is that words that weren't included in the initial vocabulary can't be assigned an embedding later on, word2vec can't assign new embeddings on the fly. With these strengths and drawbacks in mind, let's consider a more recent alternative. 


% Transformers
The paper \textit{Attention is all you need} \cite{vaswani_attention_2017} by Vaswani et al. first describes the Transformer architecture, now prominent in Large Language Models. In contrast to Word2Vec, embeddings in this architecture consist of the vector associated with the token according to the vocabulary plus a positional encoding (a combination of sine and cosine functions depending distinct for every dimension of the embedding). Additionally, although the training process bears some resemblance to the CBOW variant of Word2Vec, there are some major improvements. Instead of just predicting \textbf{a} word in a sentence based on an unordered bag consisting of previous and following words, the training process used a sequence of words as an input and the task was to predict the next word. The ''byproduct'' of this process is essentially one of the earlier versions of ChatGPT - a model that predicts what the next word (or rather token, which might also be something like a syllable) should be in a conversation repeatedly. The internal architecture to make predictions differs heavily from the initial Word2Vec model. Several components are stacked in parallel and sequentially to ensure that each word embedding incorporates the context of all of it's neighborhood (or only its predecessors, depending on the mask used). From a birds's eye view you can separate the transformer into two parts: The encoder and the decder. The encoder takes the initial input text, passes it through its network, and then passes it on to the decoder. Depending on training or inference time, the decoder either gets the target text or the generated response up until this point as an input, passes it through it's first layers, merges it together with the data received from the encoder, passes the result through the following layers and uses that to assign probabilities to all possible output-tokens to indicate their likelihood of being the next in the sequence. The internals of the two blocks each consist of multiple layers each, the most important part, which will briefly explained, are the \textit{self-attention}/\textit{multi-head-attention} layers. Self-attention lies of the core of incorporating the context and, as the paper title suggests, is all you need. Attention functions consist of a query and a set of key-value pairs as input, and their output is a weighted sum of the values biased by some compatibility function applied to the query and the corresponding key. This particular type of attention corresponds to this function: 
$$\mathbf{Attention}(Q,K,V) = softmax(\frac{Q \cdot K^T}{\sqrt{d_k}}) \cdot V$$
In a scenario where all of these are vectors or rather embeddings, it has an additional benefit of being easily parallelizable by using matrices instead, thereby enabling highly efficient GPU computations.\\
Self-attention is used in several different parts of a transformer. In the encoder block, it is used so that a singular word (or token) embedding incorporates the context of surrounding words and semantic connections become present. The paper even states that it appears to be benefitial to use several Self-attention layers in parallel, concatenating the results and applying a linear transformation, which suggests that there are several distinct levels of semantics that connect different parts of a text corpus (that are learned). This is called multi-head attention. The next use in the architecture is in the first layers of the decoder stack. The idea is the same as in the encoder, however, every word embedding is only allowed to use its \textbf{predecessors} as context, following embeddings are masked out. This is to ensure that the model learns to predict the next word in the output and doesn't use information present only in the training as an unwanted shortcut. The final use of an attention layer lies in the mergeing of encoder and decoder stacks. In this case  the question posed is ''Which parts of the input are related to which parts of the output''. Or, in terms of the paper ''This allows every position in the decoder to attend over all positions in the input sequence.'' \cite[p. 5]{vaswani_attention_2017}

% BERT
The final step in the ladder regarding Text embeddings that will be briefly explained, is BERT, which stands for Bidirectional Encoder Representations from Transformers. While the initial transformers purpose was purely generative in nature, BERT is more a jack-of-all-trades type of model. It's used for tasks such as question answering where it finds the answer to a question within a given paragraph, sentiment analysis, classifying inferences between two sentences, semantic similarities and many more. There are several new concepts it introduced: Fully bidirectional attention mechanisms without masking - while the initial Transformer architecture by Vaswani et al. suggested Masking in the self-attention layers to prevent ''cheating'', BERT only uses masks in the actual training data, there is no masking in the entire architecture. This enables word embeddings to attend to the entire context, not only predecessors or successors. The way this is achieved is by altering the (pre-)training task - instead of predicting the next token, parts of the input are randomly masked and the task of the model is to predict the initial token which was supposed to be there. Additionally, the paper demonstrated how a single pretrained BERT instance can easily be finetuned to several different tasks - and outperformed many task-specific architectures at the time. The general architecture is identical to the encoder stack of the initial transformer architecture discussed previously. An input can be a single piece of text, or two connected halves, which are tokenized by a vocabulary that is built up of likelihood-based word and subword fragments with WordPiece \cite{wu_googles_2016} and then fed through the neural network. A special addition is a starting token ($\mathtt{[CLS]}$), which is prepended to each statement. This can then used for classification tasks later on. The $\mathtt{[SEP]}$ token is used to split the input into two segments, for tasks like quesion answer pairs and sentence inference. The output of the neural network is a set of text embeddings for each token of the input with context information from all other tokens. These are what can then be used for the previously mentioned downstream tasks. The pre-training consists of the aforementioned masking, dubbed ''Masked Language Model'' or MLM for short, and Next Sentence Prediction (NSP). For NSP, the model gets two sentences and has to predict if the second one follows the first in the original text corpus, or if it is a non-connected random sentence. For this task, the embedding of the $\mathtt{[CLS]}$ token at the beginning of the sequence is used. \cite{BERT}

\section{Knowledge Graphs}
A Knowledge Graph (KG) can formally be described as a set of triples of the form 
$\mathtt{< head,\ relation,\ tail>}$. The ''Knowledge'' part from a Knowledge Graph is that each part of a triple is supposed to reference some sort of real-world entity or concept. Therefore each reference has some sort notion of semantics with additional \textit{Knowledge} attached, which can be explicitly or implicitly modeled in the Knowledge Graph. Instances of Knowledge Graphs include (but aren't limited to) the financial domain \cite{Bellomarini_Magnanimi_Nissl_Sallinger_2021}, where ownership, trade and similar relations form the links between people, companies or even entire countries, the bio-medical domain \cite{Lu_Goi_Zhao_Wang_2025}, where molecular structures, proteins etc. are used to model interactions and likelihoods of diseases, toxins and other substances relevant to modern medicine, and many more.

There are several benefits from using Knowledge Graphs to model your domain. Firstly, the aforementioned data model is very simple and can be represented in all kinds of data stores (not limited to graph stores!). Additionally, it's a powerful abstraction that enables sophisticated reasoning over a simple interface while still being able to tap into all the underlying complexity at will. Finally, there are well-established methodologies from the fields of reasoning and inference as well as machine learning, which are easily applicable to KGs. From the reasoning/database disciplines, there are tools such as Datalog, Answer Set Programming and Ontology-based reasoning such as the Semantic Web Stack with OWL and RDF. As for machine learning, models such as Knowlege Graph Embeddings (KGEs) and Graph Neural Networks (GNNs) are prime examples. GNNs will be further discussed in the next section, as they are a key part of this thesis. KGEs are a model that tries to map a set of triples (and thereby a Knowledge Graph) in a vector domain, where the nodes (which are represented by the heads and tails of the triples) are represented as points in this domain and relations are some sort of operation or formula for a pair of nodes to determine if the model predicts there to be an edge between the two or not. A classic example from literature is TransE, where relations are simple offset vectors and an edge is predicted as true iff the following equation holds true: $$embed(\mathtt{head}) + embed(\mathtt{relation}) \approx embed(\mathtt{tail})$$. 

The reason why this isn't well-suited for this thesis is that KGEs \textit{only} learn the graph structure and features of nodes (and edges) aren't incorporated into the model's embeddings.


\section{Graph Neural Networks}
In this section the general concept of a Graph Neural network (GNN) will be explored, followed by a more in-depth exploration of two specific models, Graph Convolutional Neural Networks (GCNs) andGraphSAGE.
A Graph Neural Network is a special type of a Multi-Layer Perceptron which is well established in Machine Learning. As the name implies, its intended for graph structures, which make it an important tool for to leverage Knowledge Graphs. The basic concept is that of a message passing network, where each node carries certain information (represented as a feature vector), which it's neighbors can see and have access to and let it influence their own feature vectors. Repeating this accumulation of data from the neighborhood and allowing it to bias a node's feature vector means more and more general structural information from further distances will accumulate in the individual nodes. A general equation that captures this notion of aggregating and adjusting that all GNNs share is the following:
$$h^{n+1}_i = \delta(h^n_i,AGG(\{h^n_j | v_j \in \mathcal{N}_{v_i}\}))$$
$h^n_i$ represents the feature vector of vertex $v_i$ after $n$ iterations, $\delta$ and $AGG$ are activation and aggregation functions that differ between models and $\mathcal{N}_{v_i}$ is the set of vertices adjacent to $v_i$ (with no clear order). Note that this means $AGG$ takes an \textit{unordered} sequence of nodes as an input and has to be permutation-independent function. 
Their applications range from recommender systems (which will be a focus of this thesis), over social networks, traffic prediction, financial and biological applications to Computer Vision and text classification problems. Needless to say, they are a very powerful tool with a variety of applications. \cite{Explainable_GNNs}


As the name suggests, Graph Convolutional Neural Networks leverage the generalization of convolution to graphs as the method to bias feature vectors. 

GraphSAGE is a framework to generalize GNNs to an inductive setting, where you can embed unseen nodes more easily and also embed large graphs without loading the whole adjacency matrix and node features into main memory at once. Their main innovation is the fact that their entire training process uses only batches of neighborhoods at a time - thereby enabling both aforementioned properties. 



\chapter{Practical Considerations}
The code in the repository can be split into to distinct parts. The embedding pipeline, which was the main focus, and the webserver providing the search functionality, which are both built in python. The embedding pipeline utilizes 
The main technologies used for this thesis were python with several important libraries, to provide the necessary functionality, and docker, for ease of development and deployment. The main packages used were pytorch and pytorch\_geometric, for the GNN training, sentence transformers, which is a package provided by Huggingface, an online hub for publicizing open source models and was used for obtaining the initial embedding model, and fastAPI, which is a python library meant to quickly set up a webserver.
For this thesis, a Bert model by the Deutsche Telekom company was used. This model was published to Huggingface. \cite{german_bert}

\section{Pipeline}

\section{Webserver and Deployment}


\section{Architecture and Code}
The search is built as a standalone docker container with all secrets passed via an env file. The results are written to an external volume specified in the docker-compose file, where an sqlite3 file contains the knowledge graph and all node embeddings, while a subfolder contains all files downloaded from the qibri backend and, where possible, an additional file containing the extracted content using Chunkr \footnote{An external service for extracting text from a variety of files \url{https://chunkr.ai/}}. The sqlite3 database is checked into a separate git repository to make it possible to go to a specific run of the pipeline and compare results.

There are two different entry points for the program. One is a webserver, to provide the actual search as a microservice. The second is the pipeline, which starts a fresh extraction of the Knowledge Graph of a specific customer from qibri, writes that to the SQLite file (after dropping it), downloads potentially missing files, runs the GNN training loop and finally writes the obtained embeddings to the SQLite file as well, thereby providing the vectors used by the search in the webserver.


\chapter{Results}

% Remove following line for the final thesis.
%%% intro.tex
%% Copyright (C) 2014-2024 by Thomas Auzinger <thomas@auzinger.name>
%
% This work may be distributed and/or modified under the
% conditions of the LaTeX Project Public License, either version 1.3
% of this license or (at your option) any later version.
% The latest version of this license is in
%   http://www.latex-project.org/lppl.txt
% and version 1.3 or later is part of all distributions of LaTeX
% version 2005/12/01 or later.
%
% This work has the LPPL maintenance status `maintained'.
%
% The Current Maintainer of this work is Thomas Auzinger.
%
% This work consists of the files vutinfth.dtx and vutinfth.ins
% and the derived file vutinfth.cls.
% This work also consists of the file intro.tex.


\newacronym{ctan}{CTAN}{Comprehensive TeX Archive Network}
\newacronym{faq}{FAQ}{Frequently Asked Questions}
\newacronym{pdf}{PDF}{Portable Document Format}
\newacronym{svn}{SVN}{Subversion}
\newacronym{wysiwyg}{WYSIWYG}{What You See Is What You Get}

\newglossaryentry{texteditor}
{
  name={editor},
  description={A text editor is a type of program used for editing plain text files.}
}

\chapter{Introduction to \LaTeX}

Since \LaTeX\ is widely used in academia and industry, there exists a plethora of freely accessible introductions to the language.
Reading through the guide at \url{https://en.wikibooks.org/wiki/LaTeX} serves as a comprehensive overview of most of the functionality and is highly recommended before starting with a thesis in \LaTeX.

\section{Installation}

A full \LaTeX\ distribution\index{distribution} consists not only of the binaries that convert the source files to the typeset documents but also of a wide range of packages and their documentation.
Depending on the operating system, different implementations are available as shown in Table~\ref{tab:distrib}.
\textbf{Due to the large number of packages that are in everyday use and due to their high interdependence, it is paramount to keep the installed distribution\index{distribution} up to date.}
Otherwise, obscure errors and tedious debugging ensue.

\begin{table}
  \centering
  \begin{tabular}{cccc}
    \toprule
    Distribution & Unix         & Windows      & MacOS        \\
    \midrule
    TeX Live     & \textbf{yes} & yes          & (yes)        \\
    MacTeX       & no           & no           & \textbf{yes} \\
    MikTeX       & (yes)        & \textbf{yes} & yes          \\
    \bottomrule
  \end{tabular}
  \caption{\TeX/\LaTeX\ distributions for different operating systems. Recommended choice is in \textbf{bold}.}
  \label{tab:distrib} % \label has to be placed AFTER \caption to produce correct cross-references.
\end{table}

\section{Editors}

A multitude of \TeX\ \glspl{texteditor} are available differing in their editing models, their supported operating systems, and their feature sets.
A comprehensive overview of \glspl{texteditor} can be found on the Wikipedia page \url{https://en.wikipedia.org/wiki/Comparison_of_TeX_editors}.
TeXstudio (\url{http://texstudio.sourceforge.net/}) is recommended.
Most editors support synchronization of the generated document and the \LaTeX\ source by \verb|Ctrl|-clicking either on the source document or the generated document.

\section{Compilation}

Modern editors usually provide the compilation programs to generate \gls{pdf} documents and for most \LaTeX\ source files, this is sufficient.
More advanced \LaTeX\ functionality, such as glossaries and bibliographies, needs additional compilation steps, however.
It is also possible that errors in the compilation process invalidate intermediate files and force subsequent compilation runs to fail.
It is advisable to delete intermediate files (\verb|.aux|, \verb|.bbl|, etc.), if errors occur and persist.
All files that are not generated by the user are automatically regenerated.
To compile the current document, the steps as shown in Table~\ref{tab:compile} have to be taken.


\begin{table}
  \centering
  \begin{tabular}{rl}
    \toprule
    & Description \\
    \midrule
    1 & Scan for refs, toc/lof/lot/loa items and cites \\
    2 & Build the bibliography     \\
    3 & Link refs and build the toc/lof/lot/loa \\
    4 & Link the bibliography \\
    5 & Build the glossary \\
    6 & Build the acronyms \\
    7 & Build the index \\
    8 & Link the glossary, acronyms, and the index \\
    9 & Link the bookmarks \\
    \midrule
    & Command \\
    \midrule
    1 & \verb|pdflatex.exe  example| \\
    2 & \verb|bibtex.exe    example| \\
    3 & \verb|pdflatex.exe  example| \\
    4 & \verb|pdflatex.exe  example| \\
    5 & \verb|makeindex.exe -t example.glg -s example.ist| \\
      & \verb|              -o example.gls example.glo| \\
    6 & \verb|makeindex.exe -t example.alg -s example.ist| \\
      & \verb|              -o example.acr example.acn| \\
    7 & \verb|makeindex.exe -t example.ilg -o example.ind example.idx| \\
    8 & \verb|pdflatex.exe  example| \\
    9 & \verb|pdflatex.exe  example| \\
    \bottomrule
  \end{tabular}
  \caption{Compilation steps for this document. The following abbreviations were used: table of contents (toc), list of figures (lof), list of tables (lot), list of algorithms (loa).}
  \label{tab:compile} % \label has to be placed AFTER \caption to produce correct cross-references.
\end{table}


\section{Basic Functionality}

In this section, various examples are given of the fundamental building blocks used in a thesis.
Many \LaTeX\ commands have a rich set of options that can be supplied as optional arguments.
The documentation of each command should be consulted to get an impression of the full spectrum of its functionality.

\subsection{Floats}

Two main categories of page elements can be differentiated in the usual \LaTeX\ workflow:
%
\begin{enumerate*}[label={\itshape(\roman*)}]
	\item the main stream of text and
	\item floating containers that are positioned at convenient positions throughout the document.
\end{enumerate*}
%
In most cases, tables, plots, and images are put into such containers since they are usually positioned at the top or bottom of pages.
These are realized by the two environments \verb|figure| and \verb|table|, which also provide functionality for cross-referencing (see Table~\ref{tab:intro} and Figure~\ref{fig:intro}) and the generation of corresponding entries in the list of figures and the list of tables.
Note that these environments solely act as containers and can be assigned arbitrary content.

\subsection{Tables}

A table in \LaTeX\ is created by using a \verb|tabular| environment or any of its extensions, e.g., \verb|tabularx|.
The commands \verb|\multirow| and \verb|\multicolumn| allow table elements to span multiple rows and columns.

\begin{table}[ht] % placement specifier
  \centering
  \begin{tabular}{lll}
    \toprule
    \multicolumn{2}{c}{Position} \\
    \cmidrule{1-2} % partial horizontal rule
    Group & Abbrev & Name \\
    \midrule
    Goalkeeper & GK & Paul Robinson \\
    \midrule
    \multirow{4}{*}{Defenders} & LB & Lucus Radebe \\
                               & DC & Michael Duburry \\
                               & DC & Dominic Matteo \\
                               & RB & Didier Domi \\
    \midrule
    \multirow{3}{*}{Midfielders} & MC & David Batty \\
                                 & MC & Eirik Bakke \\
                                 & MC & Jody Morris \\
    \midrule
    Forward & FW & Jamie McMaster \\
    \midrule
    \multirow{2}{*}{Strikers} & ST & Alan Smith \\
                              & ST & Mark Viduka \\
    \bottomrule
  \end{tabular}
  \caption[\LaTeX\ table example with shortened caption for the list of tables]{Adapted example from the \LaTeX\ guide at \url{https://en.wikibooks.org/wiki/LaTeX/Tables}. This example uses rules specific to the \texttt{booktabs} package and employs the multi-row functionality of the \texttt{multirow} package.}
  \label{tab:intro} % \label has to be placed AFTER \caption to produce correct cross-references.
\end{table}

\subsection{Images}

An image is added to a document via the \verb|\includegraphics| command as shown in Figure~\ref{fig:intro}.
The \verb|\subcaption| command can be used to reference subfigures, such as Figure~\ref{fig:intro:full width} and~\ref{fig:intro:half width}.

\begin{figure}[ht]
  \centering
  \begin{subfigure}[b]{0.45\columnwidth}
    \centering
    \includegraphics[width=\textwidth]{TUWI-Logo-Code.png}
    \subcaption{The TU Wien Informatics logo at text width.}
    \label{fig:intro:full width}
  \end{subfigure}
  \begin{subfigure}[b]{0.45\columnwidth}
    \centering
    \includegraphics[width=0.5\textwidth]{TUWI-Logo-Code.png}
    \subcaption{The TU Wien Informatics logo at half the text width.}
    \label{fig:intro:half width}
  \end{subfigure}
  \caption[Optional caption for the figure list (often used to abbreviate long captions)]{The header logo at different sizes.} % Remove the [...] argument if the original caption should be used in the figure list.
  \label{fig:intro} % \label has to be placed AFTER \caption (or \subcaption) to produce correct cross-references.
\end{figure}

\subsection{Mathematical Expressions}

One of the original motivations for creating the \TeX\ system was the need for mathematical typesetting.
To this day, \LaTeX\ is the preferred system to write math-heavy documents and a wide variety of functions aids the author in this task.
A mathematical expression can be inserted inline as $\sum_{n=1}^{\infty} \frac{1}{n^2} = \frac{\pi^2}{6}$ outside of the text stream as \[ \sum_{n=1}^{\infty} \frac{1}{n^2} = \frac{\pi^2}{6} \] or as a numbered equation with
\begin{equation}
\sum_{n=1}^{\infty} \frac{1}{n^2} = \frac{\pi^2}{6}.
\end{equation}

\subsection{Pseudo Code}

The presentation of algorithms can be achieved with various packages; the most popular are \verb|algorithmic|, \verb|algorithm2e|, \verb|algorithmicx|, or \verb|algpseudocode|.
An overview is given at \url{https://tex.stackexchange.com/questions/229355}.
An example of the use of the \verb|alogrithm2e| package is given with Algorithm~\ref{alg:gauss-seidel}.

\begin{algorithm}
  \SetKw{BreakFor}{break for}
  \KwIn{A scalar~$\epsilon$, a matrix $\mathbf{A} = (a_{ij})$, a vector $\vec{b}$, and an initial vector $\vec{x}^{(0)}$}
  \KwOut{$\vec{x}^{(n)}$ with $\mathbf{A} \vec{x}^{(n)} \approx \vec{b}$}
  \For{$k\leftarrow 1$ \KwTo maximum iterations}
  {
     \For{$i\leftarrow 1$ \KwTo $n$}
     {
        $x_i^{(k)} = \frac{1}{a_{ii}} \left(b_i-\sum_{j<i} a_{ij} x_j^{(k)} - \sum_{j>i} a_{ij} x_j^{(k-1)} \right)$\;
     }
     \If{$\lvert\vec{x}^{(k)}-\vec{x}^{(k-1)}\rvert < \epsilon$}
     {\BreakFor\;}
  }
  \Return{$\vec{x}^{(k)}$\;}
  \caption{Gauss-Seidel}
  \label{alg:gauss-seidel} % \label has to be placed AFTER \caption to produce correct cross-references.
\end{algorithm}

\section{Bibliography}

The referencing of prior work is a fundamental requirement of academic writing and is well supported by \LaTeX.
The \textsc{Bib}\TeX\ reference management software is the most commonly used system for this purpose.
Using the \verb|\cite| command, it is possible to reference entries in a \verb|.bib| file out of the text stream, e.g., as~\cite{Turing1936}.
The generation of the formatted bibliography needs a separate execution of \verb|bibtex.exe| (see Table~\ref{tab:compile}).

\section{Table of Contents}

The table of contents is automatically built by successive runs of the compilation, e.g., of \verb|pdflatex.exe|.
The command \verb|\setsecnumdepth| allows the specification of the depth of the table of contents and additional entries can be added to the table of contents using \verb|\addcontentsline|.
The starred versions of the sectioning commands, i.e., \verb|\chapter*|, \verb|\section*|, etc., remove the corresponding entry from the table of contents.

\section{Acronyms / Glossary / Index}

The list of acronyms, the glossary, and the index need to be built with a separate execution of \verb|makeindex| (see Table~\ref{tab:compile}).
Acronyms have to be specified with \verb|\newacronym| while glossary entries use \verb|\newglossaryentry|.
Both are then used in the document content with one of the variants of \verb|\gls|, such as \verb|\Gls|, \verb|\glspl|, or \verb|\Glspl|.
Index items are simply generated by placing \verb|\index|\marg{entry} next to all the words that correspond to the index entry \meta{entry}.
Note that many enhancements exist for these functionalities and the documentation of the \verb|makeindex| and the \verb|glossaries| packages should be consulted.

\section{Tips}

Since \TeX\ and its successors do not employ a \gls{wysiwyg} editing scheme, several guidelines improve the readability of the source content:
\begin{itemize}
\item Each sentence in the source text should start with a new line.
      This helps not only the user navigate through the text but also enables revision control systems (e.g. \gls{svn}, Git) to show the exact changes authored by different users.
      Paragraphs are separated by one (or more) empty lines.
\item Environments, which are defined by a matching pair of \verb|\begin{name}| and \verb|\end{name}|, can be indented by whitespace to show their hierarchical structure.
\item In most cases, the explicit use of whitespace (e.g. by adding \verb|\hspace{4em}| or \verb|\vspace{1.5cm}|) violates typographic guidelines and rules.
      Explicit formatting should only be employed as a last resort and, most likely, better ways to achieve the desired layout can be found by a quick web search.
\item The use of bold or italic text is generally not supported by typographic considerations and the semantically meaningful \verb|\emph{|\texttt{$\dots$}\verb|}| should be used.
\end{itemize}

The predominant application of the \LaTeX\ system is the generation of \gls{pdf} files via the \textsc{Pdf}\LaTeX\ binaries.
In the current version of \textsc{Pdf}\LaTeX, it is possible that absolute file paths and user account names are embedded in the final \gls{pdf} document.
While this poses only a minor security issue for all documents, it is highly problematic for double-blind reviews.
The process shown in Table~\ref{tab:ps2pdf} can be employed to strip all private information from the final \gls{pdf} document.

\begin{table}[ht]
  \centering
  \begin{tabular}{rl}
  \toprule
  & Command \\
  \midrule
  1 & Rename the \gls{pdf} document \verb|final.pdf| to \verb|final.ps|. \\
  2 & Execute the following command: \\
    & \verb|ps2pdf -dPDFSETTINGS#/prepress ^| \\
    & \verb| -dCompatibilityLevel#1.4 ^| \\
    & \verb| -dAutoFilterColorImages#false ^| \\
    & \verb| -dAutoFilterGrayImages#false ^| \\
    & \verb| -dColorImageFilter#/FlateEncode ^| \\
    & \verb| -dGrayImageFilter#/FlateEncode ^| \\
    & \verb| -dMonoImageFilter#/FlateEncode ^| \\
    & \verb| -dDownsampleColorImages#false ^| \\
    & \verb| -dDownsampleGrayImages#false ^| \\
    & \verb| final.ps final.pdf| \\
  \bottomrule
  \end{tabular}

  On Unix-based systems, replace \verb|#| with \verb|=| and \verb|^| with \verb|\|.
  \caption{Anonymization of \gls{pdf} documents.}
  \label{tab:ps2pdf}
\end{table}

\section{Resources}

\subsection{Useful Links}

In the following, a listing of useful web resources is given.
\begin{description}
\item[\url{https://en.wikibooks.org/wiki/LaTeX}] An extensive wiki-based guide to \LaTeX.
\item[\url{http://www.tex.ac.uk/faq}] A (huge) set of \gls{faq} about \TeX\ and \LaTeX.
\item[\url{https://tex.stackexchange.com/}] The definitive user forum for non-trivial \LaTeX-related questions and answers.
\end{description}

\subsection[Comprehensive TeX Archive Network]{\gls{ctan}}

The \gls{ctan} is the official repository for all \TeX-related material.
It can be accessed via \url{https://www.ctan.org/} and hosts (among other things) a huge variety of packages that provide extended functionality for \TeX\ and its successors.
Note that most packages contain \gls{pdf} documentation that can be directly accessed via \gls{ctan}.

In the following, a short, non-exhaustive list of relevant \gls{ctan}-hosted packages are given together with their relative path.
\begin{description}[itemsep=0ex]
\item[\href{https://www.ctan.org/pkg/algorithm2e}{algorithm2e}] Functionality for writing pseudo code.
\item[\href{https://www.ctan.org/pkg/amsmath}{amsmath}] Enhanced functionality for typesetting mathematical expressions.
\item[\href{https://www.ctan.org/pkg/amsfonts}{amssymb}] Provides a multitude of mathematical symbols.
\item[\href{https://www.ctan.org/pkg/booktabs}{booktabs}] Improved typesetting of tables.
\item[\href{https://www.ctan.org/pkg/enumitem}{enumitem}] Control over the layout of lists (\verb|itemize|, \verb|enumerate|, \verb|description|).
\item[\href{https://www.ctan.org/pkg/fontenc}{fontenc}] Determines font encoding of the output.
\item[\href{https://www.ctan.org/pkg/glossaries}{glossaries}] Create glossaries and lists of acronyms.
\item[\href{https://www.ctan.org/pkg/graphicx}{graphicx}] Insert images into the document.
\item[\href{https://www.ctan.org/pkg/inputenc}{inputenc}] Determines encoding of the input.
\item[\href{https://www.ctan.org/pkg/l2tabu}{l2tabu}] A description of bad practices when using \LaTeX.
\item[\href{https://www.ctan.org/pkg/mathtools}{mathtools}] Further extension of mathematical typesetting.
\item[\href{https://www.ctan.org/pkg/memoir}{memoir}] The document class upon which the \verb|vutinfth| document class is based.
\item[\href{https://www.ctan.org/pkg/multirow}{multirow}] Allows table elements to span several rows.
\item[\href{https://www.ctan.org/pkg/pgfplots}{pgfplots}] Function plot drawings.
\item[\href{https://www.ctan.org/pkg/pgf}{pgf/TikZ}] Creating graphics inside \LaTeX\ documents.
\item[\href{https://www.ctan.org/pkg/subcaption}{subcaption}] Allows the use of subfigures and enables their referencing.
\item[\href{https://www.ctan.org/tex-archive/info/symbols/comprehensive/}{symbols/comprehensive}] A listing of around 5000 symbols that can be used with \LaTeX.
\item[\href{https://www.ctan.org/pkg/voss-mathmode}{voss-mathmode}] A comprehensive overview of typesetting mathematics in \LaTeX.
\item[\href{https://www.ctan.org/pkg/xcolor}{xcolor}] Allows the definition and use of colors.
\end{description} % A short introduction to LaTeX.

\chapter{Discussion}


\backmatter

% Declare the use of AI tools as mentioned in the statement of originality.
% Use either the English aitools or the German kitools.
\begin{aitools}
\todo{Ihr Text hier.}
\end{aitools}

\begin{kitools}
\todo{Enter your text here.}
\end{kitools}

% Use an optional list of figures.
\listoffigures % Starred version, i.e., \listoffigures*, removes the toc entry.

% Use an optional list of tables.
\cleardoublepage % Start list of tables on the next empty right hand page.
\listoftables % Starred version, i.e., \listoftables*, removes the toc entry.

% Use an optional list of alogrithms.
%\listofalgorithms
%\addcontentsline{toc}{chapter}{List of Algorithms}

% Add an index.
\printindex

% Add a glossary.
\printglossaries

% Add a bibliography.
\bibliographystyle{alpha}
\bibliography{intro}

\end{document}
